\chapter{Directional Derivatives and Vector-Valued Functions}

The partial derivatives $f_{x}$ and $f_{y}$ measure the response of $f$
when one input moves and the others stand still. Real adjustments are
rarely so tidy: a firm re-balances capital and labour together, a consumer
slides along a budget line, a hiker walks north-east rather than due
north. The \textbf{directional derivative} measures the rate of change of
$f$ per unit distance travelled in any prescribed direction, and a single
dot product with the gradient computes it. The gradient thereby earns its
geometric role: it points where $f$ climbs fastest, its length is that
fastest rate, and it stands perpendicular to the level curves, to the
indifference curves and isoquants of the previous chapters.

The second half of the chapter lets the function itself be vector-valued.
A map $\mathbf{f}:\R^{n}\to\R^{m}$ (a production plan sending inputs to
(output, cost), a coordinate change, a demand system) is differentiated
by collecting the partial derivatives of all components into the
\textbf{Jacobian matrix}, the correct ``derivative as linear
approximation'' for vector-valued functions. In this language the chain
rule becomes matrix multiplication. Determinants, invertibility and the
Implicit Function Theorem built on top of this matrix are the subject of
the next chapter.

\section{Core Rules}

\begin{rulebox}{Directional Derivative}
Let $f:\R^{n}\to\R$ be differentiable at $\mathbf{x}_{0}$ and let
$\mathbf{u}\in\R^{n}$ be a \textbf{unit} vector, $|\mathbf{u}|=1$. The
directional derivative of $f$ at $\mathbf{x}_{0}$ in the direction
$\mathbf{u}$ is
\[
  D_{\mathbf{u}}f(\mathbf{x}_{0})
  \;=\; \lim_{h\to 0}\frac{f(\mathbf{x}_{0}+h\mathbf{u})
        -f(\mathbf{x}_{0})}{h}
  \;=\; \grad f(\mathbf{x}_{0})\cdot\mathbf{u} .
\]
Two variables: $D_{\mathbf{u}}f=f_{x}u_{1}+f_{y}u_{2}$. The partials are
the special cases $\mathbf{u}=(1,0)$ and $\mathbf{u}=(0,1)$. The unit
length is what makes $D_{\mathbf{u}}f$ a rate \emph{per unit distance}.
\end{rulebox}

\begin{rulebox}{Normalising the Direction}
If a direction is given by a nonzero vector $\mathbf{v}$ that is not of
unit length, normalise first:
\[
  \mathbf{u} \;=\; \frac{\mathbf{v}}{|\mathbf{v}|}
  \;=\; \frac{\mathbf{v}}{\sqrt{v_{1}^{2}+\cdots+v_{n}^{2}}} .
\]
``The direction from $P$ to $Q$'' means $\mathbf{v}=Q-P$, then normalise.
Forgetting this step is a frequent error in directional-derivative
problems.
\end{rulebox}

\begin{rulebox}{Steepest Ascent and Descent}
Since $D_{\mathbf{u}}f=\grad f\cdot\mathbf{u}=|\grad f|\cos\theta$, where
$\theta$ is the angle between $\mathbf{u}$ and $\grad f$, at any point
with $\grad f(\mathbf{x}_{0})\neq\mathbf{0}$:
\begin{itemize}[leftmargin=1.6em,itemsep=0.3ex]
  \item $D_{\mathbf{u}}f$ is \emph{maximal} for
        $\mathbf{u}=\grad f/|\grad f|$, with maximum value $|\grad f|$
        (steepest ascent);
  \item $D_{\mathbf{u}}f$ is \emph{minimal} for
        $\mathbf{u}=-\grad f/|\grad f|$, with value $-|\grad f|$
        (steepest descent);
  \item $D_{\mathbf{u}}f=0$ exactly when $\mathbf{u}\perp\grad f$.
\end{itemize}
For output $Q(K,L)$, the gradient $(Q_{K},Q_{L})$ is the input adjustment
that raises output fastest per unit of adjustment, and $|\grad Q|$ is that
maximal marginal gain.

\textbf{A fixed amount of adjustment.} If the change vector
$(\mathrm{d}K,\mathrm{d}L)$ must have a given length $I$, output rises
fastest when it points along the gradient:
\[
  \mathrm{d}K:\mathrm{d}L=Q_{K}:Q_{L},\qquad
  (\mathrm{d}K,\mathrm{d}L)=\frac{I}{|\grad Q|}\,(Q_{K},Q_{L}),\qquad
  \mathrm{d}Q\approx I\,|\grad Q| .
\]
A \emph{budget} $\mathrm{d}K+\mathrm{d}L=I$ with $\mathrm{d}K,\mathrm{d}L\ge0$ is
a different constraint: the linear estimate $Q_{K}\mathrm{d}K+Q_{L}\mathrm{d}L$ is then largest at a
corner, with everything spent on the input with the larger marginal
product.
\end{rulebox}

\begin{rulebox}{Gradient and Level Curves}
Along the level set $f=c$ the function does not change, so every tangent
direction $\mathbf{u}$ of the level set satisfies
$D_{\mathbf{u}}f=\grad f\cdot\mathbf{u}=0$:
\[
  \grad f \;\perp\; \text{level curve (surface) through the point.}
\]
In two variables (with $f_{y}\neq0$) the tangent direction is proportional
to $(f_{y},-f_{x})$, consistent with the implicit slope
$\dd{y}{x}=-f_{x}/f_{y}$ of the previous chapters. Economically: the
marginal-utility vector $(U_{x},U_{y})$ is normal to the indifference
curve, and $(Q_{K},Q_{L})$ is normal to the isoquant.
\end{rulebox}

\begin{rulebox}{Tangent lines: to a level curve, and to the graph}
\textbf{Tangent to the level curve} $f(x,y)=c$ through
$\mathbf{p}=(x_{0},y_{0})$, where $\grad f(\mathbf{p})\neq\mathbf{0}$. The
gradient is a normal vector, so the tangent line is
\[
  \grad f(\mathbf{p})\cdot(\mathbf{x}-\mathbf{p})=0,
  \qquad\text{i.e.}\qquad
  f_{x}(\mathbf{p})(x-x_{0})+f_{y}(\mathbf{p})(y-y_{0})=0
\]
(normal form), or $\mathbf{x}=\mathbf{p}+t\mathbf{h}$ with
$\mathbf{h}\perp\grad f(\mathbf{p})$, for instance
$\mathbf{h}=(f_{y},-f_{x})$ (parametric form). Both describe the same
line in the plane.

\textbf{Tangent to the graph in a direction.} The vertical plane through
$P=(x_{0},y_{0})$ along the unit vector $\mathbf{u}$ cuts the surface
$z=f(x,y)$ in the curve $t\mapsto f(P+t\mathbf{u})$, whose slope at $t=0$
is $D_{\mathbf{u}}f(P)$. That is the geometric meaning of the directional
derivative: the slope of the surface in the direction $\mathbf{u}$. The
tangent line to the surface in that direction is a line in $\R^{3}$:
\[
  (x,y,z)=\bigl(x_{0},\,y_{0},\,f(P)\bigr)
          +t\,\bigl(u_{1},\,u_{2},\,D_{\mathbf{u}}f(P)\bigr),\qquad t\in\R .
\]
\begin{itemize}[leftmargin=1.6em,itemsep=0.3ex]
  \item $x$-direction ($\mathbf{u}=\mathbf{e}_{1}$):
        $\bigl(x_{0}+t,\;y_{0},\;f(P)+t\,f_{x}(P)\bigr)$;
        $y$-direction ($\mathbf{u}=\mathbf{e}_{2}$):
        $\bigl(x_{0},\;y_{0}+t,\;f(P)+t\,f_{y}(P)\bigr)$.
  \item For a direction vector $\mathbf{h}$ that is not normalised,
        $\bigl(P+t\mathbf{h},\;f(P)+t\,\grad f(P)\cdot\mathbf{h}\bigr)$ is
        the \emph{same} line, run through at a different speed. Only the
        slope, rise per unit of horizontal distance, needs the unit vector.
  \item All these lines lie in the tangent plane of \chref{derivatives}{10}.
\end{itemize}
\end{rulebox}

\begin{rulebox}{Second Directional Derivative}
For a $C^{2}$ function $f$, differentiating once more in the same unit
direction $\mathbf{u}$ gives
\[
  D^{2}_{\mathbf{u}}f
  \;=\; \mathbf{u}^{T}\Hess_{f}\,\mathbf{u}
  \;=\; \sum_{i,j} f_{x_{i}x_{j}}\,u_{i}u_{j} ;
  \qquad\text{two variables:}\quad
  D^{2}_{\mathbf{u}}f
  \;=\; f_{xx}u_{1}^{2}+2f_{xy}u_{1}u_{2}+f_{yy}u_{2}^{2} .
\]
It is the second derivative at $t=0$ of the profile
$t\mapsto f(\mathbf{x}_{0}+t\mathbf{u})$: the curvature of $f$ along the
line through $\mathbf{x}_{0}$ in direction $\mathbf{u}$. Definiteness of
the Hessian (\chref{derivatives}{11}) is exactly the statement that
$D^{2}_{\mathbf{u}}f$ has one sign for \emph{every} direction
$\mathbf{u}$.
\end{rulebox}

\begin{rulebox}{Curvature in a direction: computing it, and what it proves}
\textbf{The profile method.} Without the Hessian: set
$\varphi(t)=f(\mathbf{x}_{0}+t\mathbf{u})$, a function of one variable,
and differentiate twice. Then $\varphi'(0)=D_{\mathbf{u}}f(\mathbf{x}_{0})$
and $\varphi''(0)=D^{2}_{\mathbf{u}}f(\mathbf{x}_{0})$. Doing it both ways
is a built-in check.

\textbf{Non-unit directions.} For any $\mathbf{h}\neq\mathbf{0}$, with
$\mathbf{u}=\mathbf{h}/|\mathbf{h}|$,
\[
  \mathbf{h}^{T}H\mathbf{h}=|\mathbf{h}|^{2}\,D^{2}_{\mathbf{u}}f,
  \qquad
  \grad f\cdot\mathbf{h}=|\mathbf{h}|\,D_{\mathbf{u}}f .
\]
For the \emph{sign} of the curvature, normalising is unnecessary; for the
value of $D_{\mathbf{u}}f$ or $D^{2}_{\mathbf{u}}f$, divide by $|\mathbf{h}|$
or $|\mathbf{h}|^{2}$.

\textbf{Directional verdicts.} At $\mathbf{x}_{0}$, $f$ is
\emph{convex in the direction $\mathbf{u}$} if $D^{2}_{\mathbf{u}}f>0$,
\emph{concave in that direction} if $D^{2}_{\mathbf{u}}f<0$; the profile
$\varphi$ bends up or down there. A value $0$ gives no verdict.

\textbf{Refuting global convexity.} A convex $f$ has a PSD Hessian at
\emph{every} point, so $D^{2}_{\mathbf{u}}f\ge0$ everywhere in every
direction; a concave $f$ has $D^{2}_{\mathbf{u}}f\le0$ everywhere. Hence one
point and direction with $D^{2}_{\mathbf{u}}f<0$ shows that $f$ is not
convex, and one with $D^{2}_{\mathbf{v}}f>0$ (for example at a strict local
minimum with PD Hessian) shows that $f$ is not concave. Both together prove
that $f$ is \textbf{neither}.

\textbf{Where this goes next.} The box ``Curvature at a point and in a
direction'' in \chref{optimization}{3} turns these directional verdicts into
statements about convexity near a point and on a whole domain, including
the case of a semidefinite, singular Hessian where no local conclusion is
possible.
\end{rulebox}

\begin{rulebox}{Vector-Valued Functions and the Jacobian Matrix}
For $\mathbf{f}:\R^{n}\to\R^{m}$ with components
$\mathbf{f}=(f_{1},\ldots,f_{m})$, the \textbf{Jacobian matrix} at
$\mathbf{x}$ is the $m\times n$ matrix
\[
  J_{\mathbf{f}}(\mathbf{x}) \;=\;
  \begin{pmatrix}
    \partial f_{1}/\partial x_{1} & \cdots & \partial f_{1}/\partial x_{n}\\
    \vdots & \ddots & \vdots\\
    \partial f_{m}/\partial x_{1} & \cdots & \partial f_{m}/\partial x_{n}
  \end{pmatrix}:
  \quad\text{row $i$ is the gradient of $f_{i}$.}
\]
It is the best linear approximation of $\mathbf{f}$ near $\mathbf{x}_{0}$:
\[
  \mathbf{f}(\mathbf{x}_{0}+\mathbf{h})
  \;\approx\; \mathbf{f}(\mathbf{x}_{0})
  +J_{\mathbf{f}}(\mathbf{x}_{0})\,\mathbf{h} ,
\]
the matrix version of $f(x+h)\approx f(x)+f'(x)h$. Special cases: $m=1$
gives the gradient as a $1\times n$ row; $n=1$ gives the velocity
(tangent) column vector $\mathbf{r}'(t)$ of a curve.
\end{rulebox}

\begin{rulebox}{Differentiability of vector-valued maps}
$\mathbf{f}\colon D\to\R^{m}$, $D\subseteq\R^{n}$ open, is
\textbf{differentiable} at $\mathbf{x}_{0}$ if there is an $m\times n$
matrix $J$ with
\[
  \mathbf{f}(\mathbf{x}_{0}+\mathbf{h})=\mathbf{f}(\mathbf{x}_{0})
  +J\mathbf{h}+\mathbf{r}(\mathbf{h}),
  \qquad
  \frac{\|\mathbf{r}(\mathbf{h})\|}{\|\mathbf{h}\|}\to0\ \text{ as }
  \mathbf{h}\to\mathbf{0} .
\]
The ``$\approx$'' of the Jacobian box is made exact by the remainder.
\begin{itemize}[leftmargin=1.6em,itemsep=0.3ex]
  \item The matrix $J$ is then unique and equals the Jacobian
        $J_{\mathbf{f}}(\mathbf{x}_{0})$ of partial derivatives. It is
        called the \textbf{derivative} (or total derivative,
        \de{Ableitung, Jacobi-Matrix}) of $\mathbf{f}$ at $\mathbf{x}_{0}$.
  \item $\mathbf{f}$ is differentiable exactly when every component
        $f_{i}$ is: the remainder vector tends to $0$ faster than
        $\|\mathbf{h}\|$ exactly when each of its entries does.
  \item Existence of all partials is \emph{necessary} but not sufficient
        (the function $xy/(x^{2}+y^{2})$ of \chref{derivatives}{10} is a single component that fails).
        Continuous partials on an open set are sufficient: a $C^{1}$ map is
        differentiable.
  \item Geometry: for $n=m=1$ the linear part is the tangent line; for
        $n=2$, $m=1$ it is the tangent plane; for a curve ($n=1$) it is the
        tangent vector $\mathbf{r}'(t)$.
\end{itemize}
\end{rulebox}

\begin{rulebox}{The shape of a Jacobian, and which compositions exist}
$J_{\mathbf{f}}$ has one \textbf{row per output} and one \textbf{column per
input}: (number of outputs)$\times$(number of inputs).
\begin{itemize}[leftmargin=1.6em,itemsep=0.3ex]
  \item A curve $\R\to\R^{m}$ gives an $m\times1$ column; a scalar function
        $\R^{n}\to\R$ a $1\times n$ row; $\R^{3}\to\R^{2}$ a $2\times3$
        (wide) matrix; $\R^{2}\to\R^{4}$ a $4\times2$ (tall) one.
  \item \textbf{Linear and affine maps} have constant Jacobians:
        $J(A\mathbf{x}+\mathbf{b})=A$ for an $m\times n$ matrix $A$;
        $J(c\,\mathbf{x}-d\,\mathbf{a})=c\,I_{n}$; the identity map has
        $J=I_{n}$; $J(\mathbf{a}^{T}\mathbf{x})=\mathbf{a}^{T}$.
  \item \textbf{Patterns in $n$ dimensions.} If component $k$ depends on
        $x_{k}$ alone, $f_{k}=g_{k}(x_{k})$, the Jacobian is diagonal,
        $\operatorname{diag}\bigl(g_{1}'(x_{1}),\ldots,g_{n}'(x_{n})\bigr)$.
\end{itemize}
\textbf{Which composition exists.} $\mathbf{g}\circ\mathbf{f}$ means
``first $\mathbf{f}$, then $\mathbf{g}$''; it is defined when the number of
outputs of $\mathbf{f}$ equals the number of inputs of $\mathbf{g}$. With
$\mathbf{f}\colon\R^{n}\to\R^{m}$ and $\mathbf{g}\colon\R^{m}\to\R^{p}$ the
chain rule multiplies $(p\times m)(m\times n)$; the inner sizes must
match, which is the same condition. Check it before computing anything:
often only one of $\mathbf{g}\circ\mathbf{f}$ and $\mathbf{f}\circ\mathbf{g}$
makes sense.
\end{rulebox}

\begin{rulebox}{Chain Rule in Matrix Form}
If $\mathbf{f}:\R^{n}\to\R^{m}$ and $\mathbf{g}:\R^{m}\to\R^{p}$ are
differentiable, then $\mathbf{h}=\mathbf{g}\circ\mathbf{f}$ is
differentiable and
\[
  J_{\mathbf{h}}(\mathbf{x})
  \;=\; J_{\mathbf{g}}\bigl(\mathbf{f}(\mathbf{x})\bigr)\,
        J_{\mathbf{f}}(\mathbf{x}),
  \qquad
  (p\times n) \;=\; (p\times m)(m\times n) .
\]
Evaluate the \emph{outer} Jacobian at the \emph{inner value}
$\mathbf{f}(\mathbf{x})$; the order of the factors is fixed by the matrix
shapes. Every chain rule of \chref{derivatives}{12} is the case $p=1$: for
$z=f(\mathbf{r}(t))$,
\[
  \dd{z}{t} \;=\; \grad f\bigl(\mathbf{r}(t)\bigr)\cdot\mathbf{r}'(t),
\]
a $1\times m$ row times an $m\times 1$ column.
\end{rulebox}

\section{Worked Example}

\begin{worked}{For $f(x,y)=x^{2}+3xy-y^{2}$ at $P=(1,2)$: the derivative in the direction $(3,4)$, the steepest ascent, a direction of no change; then the Jacobian of $F(x,y)=(x^{2}-y,\;xy)$}
\wstep{Step 1. The gradient first --- every question below is an inner product with it.}
\[
  f_{x}=2x+3y,\qquad f_{y}=3x-2y,
\]
\[
  \grad f(1,2)=\bigl(2+6,\;3-4\bigr)=(8,-1).
\]

\wstep{Step 2. Normalise the direction, which the formula requires.}
\[
  \|(3,4)\|=5,
  \qquad
  \mathbf u=\Bigl(\tfrac35,\tfrac45\Bigr).
\]
\[
  D_{\mathbf u}f(P)=\grad f\cdot\mathbf u=8\cdot\tfrac35+(-1)\cdot\tfrac45
  =\frac{24-4}{5}=\boxed{4}.
\]
\wtrap{Without normalising you would get $\grad f\cdot(3,4)=20$, five times
too large. The directional derivative is a rate \emph{per unit distance
travelled}, so doubling the length of the direction vector must not double
the answer. If the problem hands you a non-unit vector --- and it usually
does --- divide by its length first.}

\wstep{Step 3. Steepest ascent and descent.}
The steepest increase is along the gradient itself:
\[
  \text{direction }\frac{(8,-1)}{\sqrt{65}},
  \qquad
  \text{rate }\|\grad f\|=\sqrt{65}\approx8.06 .
\]
The steepest \emph{decrease} is the opposite direction, at rate
$-\sqrt{65}$. Every other direction lies strictly between, since
\[
  D_{\mathbf u}f=\grad f\cdot\mathbf u=\|\grad f\|\cos\theta ,
\]
and $\cos\theta\in[-1,1]$. Our $4$ sits comfortably inside
$[-8.06,\,8.06]$. \checkmark
\wtrap{The value $8.06$ is the \emph{rate}; the direction is the normalised
gradient $(0.992,-0.124)$, not $(8,-1)$. Reporting the unnormalised gradient
as ``the direction'' conflates a direction with a magnitude.}

\wstep{Step 4. A direction of no change --- and what it is tangent to.}
$D_{\mathbf u}f=0$ requires $\mathbf u\perp\grad f$. Rotating $(8,-1)$ by a
right angle gives $(1,8)$, so
\[
  \mathbf u=\frac{(1,8)}{\sqrt{65}}
  \quad\Longrightarrow\quad
  D_{\mathbf u}f=\frac{8\cdot1+(-1)\cdot8}{\sqrt{65}}=0 . \checkmark
\]
Moving in that direction keeps $f$ constant to first order, so $(1,8)$ is
tangent to the level curve of $f$ through $P$ --- which is the geometric
content of ``the gradient is perpendicular to the level curve''.

\wstep{Step 5. A vector-valued function needs a matrix of derivatives.}
$F(x,y)=(x^{2}-y,\ xy)$ has two components, each with two partials. Collect
them with \textbf{one row per component function}:
\[
  J_{F}(x,y)=
  \begin{pmatrix}
    \pd{F_{1}}{x} & \pd{F_{1}}{y}\\[0.6ex]
    \pd{F_{2}}{x} & \pd{F_{2}}{y}
  \end{pmatrix}
  =\begin{pmatrix}2x&-1\\ y&x\end{pmatrix}.
\]
\wtrap{Rows are components, columns are variables. Transposing this is the
standard error, and it matters as soon as the Jacobian is multiplied by
anything. A check: for a scalar function ($m=1$) the Jacobian is a single
\emph{row}, and the gradient is conventionally a column --- so
$J=\grad^{\!\top}$.}

\wstep{Step 6. Evaluate, and read the determinant.}
\[
  J_{F}(1,2)=\begin{pmatrix}2&-1\\2&1\end{pmatrix},
  \qquad
  \det J_{F}(1,2)=2+2=4\neq0 .
\]
A non-zero Jacobian determinant says the map is locally invertible near
$(1,2)$: the system
\[
  x^{2}-y=u,\qquad xy=v
\]
can be solved for $(x,y)$ in terms of $(u,v)$ near that point. That is the
Implicit Function Theorem of \chref{derivatives}{14} in its multivariate form, and
$\det J$ is the condition it turns on.

\wstep{Step 7. Why the Jacobian is the right generalisation.}
\begin{center}\footnotesize
\begin{tabular}{lll}
  \hline
  $f:\R\to\R$ & derivative $f'$ & a number\\
  $f:\R^{n}\to\R$ & gradient $\grad f$ & a vector\\
  $F:\R^{n}\to\R^{m}$ & Jacobian $J_{F}$ & an $m\times n$ matrix\\
  \hline
\end{tabular}
\end{center}
In each row the object is the best \emph{linear} approximation at a point,
and the chain rule is matrix multiplication:
$J_{G\circ F}=J_{G}\cdot J_{F}$, with the order fixed by the order of the
composition. The one-variable chain rule $g'(f(x))f'(x)$ is the
$1\times1$ case.
\end{worked}

\section{Practice Problems}

\begin{problems}
  \item Find $D_{\mathbf{u}}f(1,1)$ for $f(x,y)=x^{2}y$ in the direction
        $\mathbf{u}=(1/\sqrt{2},\,1/\sqrt{2})$.
  \item For $f(x,y)=\ln(x^{2}+y^{2})$ at $(1,1)$, find the directional
        derivative toward $\mathbf{v}=(2,1)$.
  \item For $f(x,y)=x^{3}-3xy^{2}$ at $(2,1)$, find the directional
        derivative in the direction of the vector from $(2,1)$ to
        $(5,5)$.
  \item For $f(x,y,z)=xyz$ at $(1,2,3)$, find $D_{\mathbf{u}}f$ where
        $\mathbf{u}=\tfrac{1}{3}(1,2,2)$.
  \item For $f(x,y)=\mathrm{e}^{x+y}$ at $(0,0)$, find $D_{\mathbf{u}}f$
        where $\mathbf{u}$ makes a $60^{\circ}$ angle with the positive
        $x$-axis.
  \item For $f(x,y)=\sqrt{x^{2}+y^{2}}$ at $(3,4)$, find the directional
        derivative in the direction $(1,0)$.
  \item A consumer with utility $U(x,y)=xy$ currently consumes
        $(x,y)=(2,3)$. Her basket shifts in the direction
        $\mathbf{v}=(3,-4)$, more $x$, less $y$. Find the rate of
        change of utility per unit distance moved and interpret the sign.
  \item For $f(x,y)=x^{2}+xy+y^{2}$ at $(1,2)$, find the direction of
        steepest ascent and the maximum rate of increase.
  \item For $f(x,y)=x^{2}+2y^{2}$ at $(1,1)$, find the direction in which
        $f$ decreases fastest and the corresponding rate.
  \item Find all unit vectors $\mathbf{u}$ along which
        $f(x,y)=x^{2}-y^{2}$ has zero rate of change at $(2,1)$.
  \item For $f(x,y,z)=x^{2}+y^{2}-z^{2}$ at $(1,2,2)$, determine the
        maximum rate of increase and the direction in which it is
        attained.
  \item A hiker stands at $(1,1)$ on the hill $z=4-x^{2}-y^{2}$. In which
        direction should she walk to descend most quickly, and how steep
        is that descent?
  \item Output is $Q(K,L)=K^{1/2}L^{1/2}$ with current input mix
        $(K,L)=(4,9)$. Find the rate of change of output in the balanced
        expansion direction $\mathbf{u}=(1,1)/\sqrt{2}$; then find the
        input direction that raises output fastest and the corresponding
        maximal rate.
  \item Show that at every point of a level curve $f(x,y)=c$ the gradient
        $\grad f$ is orthogonal to the tangent direction of the curve.
        Conclude that the marginal-utility vector $(U_{x},U_{y})$ is
        normal to the indifference curve.
  \item For $f(x,y)=x^{2}y$ at $(1,2)$ with $\mathbf{u}=(3/5,\,4/5)$,
        compute both $D_{\mathbf{u}}f$ and the second directional
        derivative
        $D^{2}_{\mathbf{u}}f=\mathbf{u}^{T}\Hess_{f}\,\mathbf{u}$.
  \item For $f(x,y)=x^{3}+y^{3}-3xy$: show that $(1,1)$ is a critical
        point, then compute $D^{2}_{\mathbf{u}}f(1,1)$ for
        $\mathbf{u}=(1,1)/\sqrt{2}$ and for
        $\mathbf{v}=(1,-1)/\sqrt{2}$. What do the signs suggest?
  \item For $f(x,y)=\mathrm{e}^{x}\cos y$ at $(0,0)$, find all unit
        vectors $\mathbf{u}$ with $D^{2}_{\mathbf{u}}f=0$.
  \item Compute the Jacobian of $\mathbf{f}(x,y)=(x+y,\;x-y)$ and its
        determinant. What is special about the Jacobian of a linear map?
  \item Compute the Jacobian of
        $\mathbf{f}(x,y)=(x^{2}y,\;x+y^{2},\;xy)$, state its shape, and
        evaluate it at $(1,2)$.
  \item Compute the Jacobian of
        $\mathbf{f}(x,y,z)=(x+y,\;y+z,\;x+z)$ and its determinant.
  \item For $\mathbf{f}(x,y)=(x^{2}-y,\;xy)$: compute $J_{\mathbf{f}}$,
        then use the linear approximation
        $\mathbf{f}(\mathbf{x}_{0}+\mathbf{h})\approx
        \mathbf{f}(\mathbf{x}_{0})+J_{\mathbf{f}}(\mathbf{x}_{0})\mathbf{h}$
        at $\mathbf{x}_{0}=(1,1)$ to estimate $\mathbf{f}(1.1,\,0.9)$,
        and compare with the exact value.
  \item Demand and supply are $D(p,q)=100-p+0.5q$ and
        $S(p,q)=-10+2p+q$, and the two equilibrium conditions are
        $G_{1}=D-S=0$ and $G_{2}=p+q-20=0$. Build the Jacobian matrix
        $\partial(G_{1},G_{2})/\partial(p,q)$ and compute its
        determinant.
  \item Let $\mathbf{g}(x,y)=(x+y,\;x-y)$ and
        $\mathbf{F}(u,v)=(u^{2},\;uv)$. Compute the Jacobian of
        $\mathbf{F}\circ\mathbf{g}$ by the matrix chain rule
        $J_{\mathbf{F}\circ\mathbf{g}}
        =J_{\mathbf{F}}(\mathbf{g}(x,y))\,J_{\mathbf{g}}(x,y)$, and
        verify by composing first.
  \item For $f(x,y)=xy^{2}$ and the curve
        $\mathbf{r}(t)=(\mathrm{e}^{t},\,t)$, compute $\dd{z}{t}$ for
        $z=f(\mathbf{r}(t))$ as the matrix product
        $J_{f}(\mathbf{r}(t))\,J_{\mathbf{r}}(t)$ (a $1\times2$ row
        times a $2\times1$ column) and verify directly.
  \item A firm converts inputs into (output, cost) via
        $\mathbf{g}(K,L)=\bigl(K^{1/2}L^{1/2},\;2K+L\bigr)$ and earns
        profit $\pi(Q,C)=10Q-C$. Use the matrix chain rule to find the
        marginal profits $(\pi_{K},\pi_{L})$ at $(K,L)=(1,1)$, and verify
        by substituting first.
\end{problems}

\section{Solutions}

\begin{solutions}
  \item $\grad f=(2xy,\,x^{2})$. At $(1,1)$: $\grad f=(2,1)$, and
        $\mathbf{u}$ is already a unit vector.
        \[
          D_{\mathbf{u}}f=(2)\frac{1}{\sqrt{2}}+(1)\frac{1}{\sqrt{2}}
            =\frac{3}{\sqrt{2}}=\frac{3\sqrt{2}}{2} .
        \]

  \item $\grad f=\Bigl(\dfrac{2x}{x^{2}+y^{2}},\,
        \dfrac{2y}{x^{2}+y^{2}}\Bigr)$. At $(1,1)$: $\grad f=(1,1)$.
        Normalise: $|\mathbf{v}|=\sqrt{4+1}=\sqrt{5}$, so
        $\mathbf{u}=(2,1)/\sqrt{5}$.
        \[
          D_{\mathbf{u}}f=\frac{1\cdot 2+1\cdot 1}{\sqrt{5}}
            =\frac{3}{\sqrt{5}}=\frac{3\sqrt{5}}{5} .
        \]

  \item Direction vector $\mathbf{v}=(5,5)-(2,1)=(3,4)$,
        $|\mathbf{v}|=5$, so $\mathbf{u}=(3/5,\,4/5)$.
        $\grad f=(3x^{2}-3y^{2},\,-6xy)$; at $(2,1)$:
        $\grad f=(12-3,\,-12)=(9,-12)$.
        \[
          D_{\mathbf{u}}f=9\cdot\frac{3}{5}+(-12)\cdot\frac{4}{5}
            =\frac{27}{5}-\frac{48}{5}=-\frac{21}{5} .
        \]

  \item $\grad f=(yz,\,xz,\,xy)$. At $(1,2,3)$: $\grad f=(6,3,2)$. The
        vector $\mathbf{u}=\tfrac13(1,2,2)$ is unit since
        $\sqrt{1+4+4}/3=1$.
        \[
          D_{\mathbf{u}}f=\frac{6\cdot 1+3\cdot 2+2\cdot 2}{3}
            =\frac{6+6+4}{3}=\frac{16}{3} .
        \]

  \item $\mathbf{u}=(\cos 60^{\circ},\,\sin 60^{\circ})
        =(1/2,\,\sqrt{3}/2)$, a unit vector.
        $\grad f=(\mathrm{e}^{x+y},\,\mathrm{e}^{x+y})$; at $(0,0)$:
        $\grad f=(1,1)$.
        \[
          D_{\mathbf{u}}f=\frac{1}{2}+\frac{\sqrt{3}}{2}
            =\frac{1+\sqrt{3}}{2} .
        \]

  \item $\grad f=\Bigl(\dfrac{x}{\sqrt{x^{2}+y^{2}}},\,
        \dfrac{y}{\sqrt{x^{2}+y^{2}}}\Bigr)$. At $(3,4)$:
        $\grad f=(3/5,\,4/5)$. With $\mathbf{u}=(1,0)$ (already unit):
        \[
          D_{\mathbf{u}}f=\frac{3}{5} .
        \]

  \item $\grad U=(y,\,x)$; at $(2,3)$: $\grad U=(3,2)$. Normalise:
        $|\mathbf{v}|=\sqrt{9+16}=5$, $\mathbf{u}=(3/5,\,-4/5)$.
        \[
          D_{\mathbf{u}}U=3\cdot\frac{3}{5}+2\cdot\Bigl(-\frac{4}{5}\Bigr)
            =\frac{9}{5}-\frac{8}{5}=\frac{1}{5} .
        \]
        The rate is positive: the marginal-utility gain from the extra
        $x$ slightly outweighs the loss from the reduced $y$, so utility
        rises, at $1/5$ per unit distance moved.

  \item $\grad f=(2x+y,\,x+2y)$. At $(1,2)$: $\grad f=(4,5)$,
        $|\grad f|=\sqrt{16+25}=\sqrt{41}$.
        \[
          \mathbf{u}=\frac{(4,5)}{\sqrt{41}},\qquad
          \text{maximum rate}=\sqrt{41} .
        \]

  \item $\grad f=(2x,\,4y)$. At $(1,1)$: $\grad f=(2,4)$,
        $|\grad f|=\sqrt{4+16}=2\sqrt{5}$. Steepest descent:
        \[
          \mathbf{u}=-\frac{(2,4)}{2\sqrt{5}}=-\frac{(1,2)}{\sqrt{5}},
          \qquad \text{rate}=-2\sqrt{5} .
        \]

  \item $\grad f=(2x,\,-2y)$. At $(2,1)$: $\grad f=(4,-2)$. We need
        $\mathbf{u}\perp(4,-2)$: $4u_{1}-2u_{2}=0$, i.e.
        $u_{2}=2u_{1}$. Normalising $(1,2)$:
        \[
          \mathbf{u}=\pm\frac{(1,2)}{\sqrt{5}} .
        \]
        \textbf{Check.} $(4,-2)\cdot(1,2)=4-4=0$. \checkmark

  \item $\grad f=(2x,\,2y,\,-2z)$. At $(1,2,2)$: $\grad f=(2,4,-4)$,
        $|\grad f|=\sqrt{4+16+16}=6$. Maximum rate $6$, attained in the
        direction
        \[
          \mathbf{u}=\frac{(2,4,-4)}{6}=\frac{(1,2,-2)}{3} .
        \]

  \item $\grad z=(-2x,\,-2y)$; at $(1,1)$: $\grad z=(-2,-2)$,
        $|\grad z|=2\sqrt{2}$. Steepest descent direction:
        \[
          \mathbf{u}=-\frac{(-2,-2)}{2\sqrt{2}}=\frac{(1,1)}{\sqrt{2}} ,
        \]
        i.e. north-east, directly away from the summit above the origin.
        The slope of that descent is $-|\grad z|=-2\sqrt{2}$.

  \item $Q_{K}=\tfrac12K^{-1/2}L^{1/2}$, $Q_{L}=\tfrac12K^{1/2}L^{-1/2}$.
        At $(4,9)$:
        \[
          Q_{K}=\frac{1}{2}\cdot\frac{1}{2}\cdot 3=\frac{3}{4},\qquad
          Q_{L}=\frac{1}{2}\cdot 2\cdot\frac{1}{3}=\frac{1}{3},
        \]
        so $\grad Q=(3/4,\,1/3)$. Balanced direction
        $\mathbf{u}=(1,1)/\sqrt{2}$:
        \[
          D_{\mathbf{u}}Q=\frac{3/4+1/3}{\sqrt{2}}
            =\frac{13/12}{\sqrt{2}}=\frac{13\sqrt{2}}{24} .
        \]
        Fastest increase along $\grad Q/|\grad Q|$ with
        \[
          |\grad Q|=\sqrt{\frac{9}{16}+\frac{1}{9}}
            =\sqrt{\frac{81+16}{144}}=\frac{\sqrt{97}}{12},\qquad
          \mathbf{u}^{*}=\frac{12}{\sqrt{97}}\Bigl(\frac34,\frac13\Bigr)
            =\frac{(9,4)}{\sqrt{97}} .
        \]
        The maximal rate is $\sqrt{97}/12$ (larger than $13\sqrt{2}/24$,
        as it must be).

  \item Let $\mathbf{r}(t)=(x(t),y(t))$ be a differentiable curve lying
        in the level set, $f(\mathbf{r}(t))=c$ for all $t$.
        Differentiate with the chain rule:
        \[
          0=\dd{}{t}\,c=\grad f(\mathbf{r}(t))\cdot\mathbf{r}'(t) .
        \]
        Hence the tangent vector $\mathbf{r}'(t)$ is orthogonal to
        $\grad f$ at every point of the curve. For $U(x,y)=c$ an
        indifference curve, this says exactly that
        $(U_{x},U_{y})$ is normal to it.

  \item First order: $\grad f=(2xy,\,x^{2})$; at $(1,2)$:
        $\grad f=(4,1)$.
        \[
          D_{\mathbf{u}}f=4\cdot\frac{3}{5}+1\cdot\frac{4}{5}
            =\frac{12}{5}+\frac{4}{5}=\frac{16}{5} .
        \]
        Second order: $f_{xx}=2y=4$, $f_{xy}=2x=2$, $f_{yy}=0$ at
        $(1,2)$.
        \[
          D^{2}_{\mathbf{u}}f
            =4\Bigl(\frac{3}{5}\Bigr)^{2}
             +2\cdot 2\cdot\frac{3}{5}\cdot\frac{4}{5}
             +0\cdot\Bigl(\frac{4}{5}\Bigr)^{2}
            =\frac{36}{25}+\frac{48}{25}=\frac{84}{25} .
        \]

  \item $\grad f=(3x^{2}-3y,\,3y^{2}-3x)$; at $(1,1)$:
        $\grad f=(0,0)$, so $(1,1)$ is a critical point. Second
        partials: $f_{xx}=6x=6$, $f_{yy}=6y=6$, $f_{xy}=-3$, so
        $\Hess_{f}(1,1)=\begin{pmatrix}6&-3\\-3&6\end{pmatrix}$. For
        $\mathbf{u}=(1,1)/\sqrt{2}$:
        \[
          D^{2}_{\mathbf{u}}f=6\cdot\frac12
            +2(-3)\cdot\frac12+6\cdot\frac12=3-3+3=3 .
        \]
        For $\mathbf{v}=(1,-1)/\sqrt{2}$:
        \[
          D^{2}_{\mathbf{v}}f=6\cdot\frac12
            +2(-3)\cdot\Bigl(-\frac12\Bigr)+6\cdot\frac12=3+3+3=9 .
        \]
        Both are positive, the surface curves upward in both
        directions. Indeed $D_{1}=6>0$ and
        $D_{2}=36-9=27>0$, so the Hessian is positive definite and
        $(1,1)$ is a strict local minimum (\chref{derivatives}{11}).

  \item $f_{xx}=\mathrm{e}^{x}\cos y$, $f_{xy}=-\mathrm{e}^{x}\sin y$,
        $f_{yy}=-\mathrm{e}^{x}\cos y$. At $(0,0)$: $f_{xx}=1$,
        $f_{xy}=0$, $f_{yy}=-1$. Writing
        $\mathbf{u}=(\cos\theta,\,\sin\theta)$:
        \[
          D^{2}_{\mathbf{u}}f=\cos^{2}\theta-\sin^{2}\theta=\cos 2\theta .
        \]
        This vanishes for $2\theta=\tfrac{\pi}{2}+k\pi$, i.e.
        $\theta=\tfrac{\pi}{4}+k\tfrac{\pi}{2}$:
        \[
          \mathbf{u}=\pm\frac{(1,1)}{\sqrt{2}}
          \quad\text{and}\quad
          \mathbf{u}=\pm\frac{(1,-1)}{\sqrt{2}} .
        \]

  \item
        \[
          J_{\mathbf{f}}=\begin{pmatrix}1&1\\1&-1\end{pmatrix},\qquad
          \det J_{\mathbf{f}}=1\cdot(-1)-1\cdot 1=-2 .
        \]
        The map is linear, $\mathbf{f}(\mathbf{x})=A\mathbf{x}$ with
        $A=\bigl(\begin{smallmatrix}1&1\\1&-1\end{smallmatrix}\bigr)$,
        and for a linear map $f_{i}=\sum_{j}A_{ij}x_{j}$ gives
        $\partial f_{i}/\partial x_{j}=A_{ij}$: the Jacobian is the
        constant matrix $A$ itself: a linear map is its own linear
        approximation.

  \item With $f_{1}=x^{2}y$, $f_{2}=x+y^{2}$, $f_{3}=xy$:
        \[
          J_{\mathbf{f}}=\begin{pmatrix}2xy&x^{2}\\1&2y\\y&x\end{pmatrix}
          \;\in\;\R^{3\times 2} .
        \]
        A $3\times 2$ matrix: one row per component, one column per
        variable; being non-square, it has no determinant. At $(1,2)$:
        \[
          J_{\mathbf{f}}(1,2)=\begin{pmatrix}4&1\\1&4\\2&1\end{pmatrix} .
        \]

  \item
        \[
          J_{\mathbf{f}}=\begin{pmatrix}1&1&0\\0&1&1\\1&0&1\end{pmatrix} .
        \]
        Expanding along the first row:
        \[
          \det J_{\mathbf{f}}
            =1\,(1\cdot 1-1\cdot 0)-1\,(0\cdot 1-1\cdot 1)+0
            =1+1=2 .
        \]

  \item $J_{\mathbf{f}}=\begin{pmatrix}2x&-1\\y&x\end{pmatrix}$. At
        $\mathbf{x}_{0}=(1,1)$:
        \[
          \mathbf{f}(1,1)=(0,\,1),\qquad
          J_{\mathbf{f}}(1,1)=\begin{pmatrix}2&-1\\1&1\end{pmatrix} .
        \]
        With $\mathbf{h}=(0.1,\,-0.1)$:
        \[
          J_{\mathbf{f}}(1,1)\,\mathbf{h}
            =\begin{pmatrix}2(0.1)+(-1)(-0.1)\\
              1(0.1)+1(-0.1)\end{pmatrix}
            =\begin{pmatrix}0.3\\0\end{pmatrix},
        \]
        so $\mathbf{f}(1.1,\,0.9)\approx(0,1)+(0.3,\,0)=(0.3,\,1)$.
        Exact value:
        \[
          \mathbf{f}(1.1,\,0.9)=(1.21-0.9,\;1.1\cdot 0.9)
            =(0.31,\;0.99) .
        \]
        The approximation is off by only $0.01$ in each component.

  \item Excess demand:
        \[
          G_{1}=D-S=(100-p+0.5q)-(-10+2p+q)=110-3p-0.5q ,
        \]
        and $G_{2}=p+q-20$. Hence
        \[
          \frac{\partial(G_{1},G_{2})}{\partial(p,q)}
            =\begin{pmatrix}-3&-0.5\\1&1\end{pmatrix},\qquad
          \det=-3\cdot 1-(-0.5)\cdot 1=-2.5 .
        \]
        The determinant is nonzero: exactly the condition the next
        chapter's Implicit Function Theorem needs to solve the system
        locally for $(p,q)$ when a parameter shifts.

  \item Inner and outer Jacobians:
        \[
          J_{\mathbf{g}}=\begin{pmatrix}1&1\\1&-1\end{pmatrix},\qquad
          J_{\mathbf{F}}(u,v)=\begin{pmatrix}2u&0\\v&u\end{pmatrix} .
        \]
        Evaluate the outer at the inner value
        $(u,v)=\mathbf{g}(x,y)=(x+y,\,x-y)$ and multiply:
        \[
          J_{\mathbf{F}\circ\mathbf{g}}
            =\begin{pmatrix}2(x+y)&0\\x-y&x+y\end{pmatrix}
             \begin{pmatrix}1&1\\1&-1\end{pmatrix}
            =\begin{pmatrix}2(x+y)&2(x+y)\\
              (x-y)+(x+y)&(x-y)-(x+y)\end{pmatrix}
        \]
        \[
          \phantom{J_{\mathbf{F}\circ\mathbf{g}}}
            =\begin{pmatrix}2(x+y)&2(x+y)\\2x&-2y\end{pmatrix} .
        \]
        \textbf{Check.} Composing first:
        $(\mathbf{F}\circ\mathbf{g})(x,y)
        =\bigl((x+y)^{2},\,(x+y)(x-y)\bigr)
        =\bigl((x+y)^{2},\,x^{2}-y^{2}\bigr)$, whose Jacobian is
        \[
          \begin{pmatrix}2(x+y)&2(x+y)\\2x&-2y\end{pmatrix} .
          \quad\checkmark
        \]

  \item $J_{f}=\bigl(y^{2}\;\;2xy\bigr)$ (a row) and
        $J_{\mathbf{r}}(t)=\begin{pmatrix}\mathrm{e}^{t}\\1\end{pmatrix}$
        (a column). At $\mathbf{r}(t)=(\mathrm{e}^{t},\,t)$ the row is
        $\bigl(t^{2}\;\;2t\mathrm{e}^{t}\bigr)$, so
        \[
          \dd{z}{t}
            =\bigl(t^{2}\;\;2t\mathrm{e}^{t}\bigr)
             \begin{pmatrix}\mathrm{e}^{t}\\1\end{pmatrix}
            =t^{2}\mathrm{e}^{t}+2t\mathrm{e}^{t}
            =\mathrm{e}^{t}(t^{2}+2t) .
        \]
        \textbf{Check.} $z=\mathrm{e}^{t}\,t^{2}$, so
        $\dd{z}{t}=\mathrm{e}^{t}t^{2}+2t\mathrm{e}^{t}$. \checkmark

  \item Inner Jacobian:
        \[
          J_{\mathbf{g}}(K,L)=\begin{pmatrix}
            \tfrac12K^{-1/2}L^{1/2}&\tfrac12K^{1/2}L^{-1/2}\\
            2&1\end{pmatrix},\qquad
          J_{\mathbf{g}}(1,1)=\begin{pmatrix}1/2&1/2\\2&1\end{pmatrix} .
        \]
        Outer Jacobian (a row, constant since $\pi$ is linear):
        $J_{\pi}=(10\;\;-1)$. Matrix chain rule:
        \[
          (\pi_{K},\,\pi_{L})
            =(10\;\;-1)\begin{pmatrix}1/2&1/2\\2&1\end{pmatrix}
            =\bigl(10\cdot\tfrac12-1\cdot 2,\;
              10\cdot\tfrac12-1\cdot 1\bigr)
            =(3,\;4) .
        \]
        \textbf{Check.} Substituting first,
        $\pi(K,L)=10K^{1/2}L^{1/2}-2K-L$, so
        $\pi_{K}=5K^{-1/2}L^{1/2}-2=5-2=3$ and
        $\pi_{L}=5K^{1/2}L^{-1/2}-1=5-1=4$ at $(1,1)$. \checkmark
        At the current input mix an extra unit of capital raises profit
        by about $3$, an extra unit of labour by about $4$.
\end{solutions}

\section{Additional Problems}

\begin{problems}
  \item For $f(x,y)=x^{2}-y^{2}+3xy$ at $(1,2)$, find the directional
        derivative toward $\mathbf{v}=(4,3)$.
  \item For $f(x,y,z)=x^{2}+2y^{2}+3z^{2}$ at $(1,1,1)$, find the
        directional derivative toward $\mathbf{v}=(1,-1,1)$.
  \item For $U(x,y)=\ln x+2\ln y$ at $(1,2)$, find the direction of
        steepest ascent and the maximum rate of increase.
  \item Find all unit vectors $\mathbf{u}$ with $D_{\mathbf{u}}f=0$ for
        $f(x,y)=xy$ at $(2,1)$.
  \item At a point $P$ the directional derivatives of $f$ toward $(1,0)$
        and toward $(0,1)$ are $3$ and $-4$. Find $\grad f(P)$, the
        maximum rate of increase at $P$, and $D_{\mathbf{u}}f$ toward
        $\mathbf{u}=(3/5,\,4/5)$.
  \item For $f(x,y)=\ln(x^{2}+y^{2})$ at $(1,0)$ with
        $\mathbf{u}=(0,1)$, compute $D_{\mathbf{u}}f$ and
        $D^{2}_{\mathbf{u}}f$, and interpret the two results.
  \item Compute the Jacobian of
        $\mathbf{f}(x,y)=(x^{2}-y^{2},\;3xy+y^{2})$ and evaluate it at
        $(1,1)$.
  \item For $f(x,y)=xy$ and $\mathbf{r}(t)=(t^{2},\,t^{3})$, compute
        $\dd{z}{t}$ for $z=f(\mathbf{r}(t))$ by the matrix chain rule
        and verify directly.
  \item Let $\mathbf{g}(s,t)=(st,\;s+t)$ and
        $\mathbf{F}(x,y)=(x+y,\;xy)$. Compute
        $J_{\mathbf{F}\circ\mathbf{g}}$ by the matrix chain rule and
        verify by composing first.
  \item Inputs produce two intermediate goods $x_{1}=2K+L$ and
        $x_{2}=K+3L$, and final output is $y=x_{1}x_{2}$. Use the matrix
        chain rule to find $\pd{y}{K}$ and $\pd{y}{L}$ at $(K,L)=(1,1)$,
        and verify by expanding $y$ as a polynomial in $K$ and $L$.
\end{problems}

Problems 11--14: tangent lines to the graph of $f$ in a given direction.

\begin{problems}[resume]
  \item For $f(x,y)=x^{2}y+y^{3}$ at $P=(1,2)$, give the slope of the graph
        and the tangent line to the graph in the $x$-direction, in the
        $y$-direction, and in the direction of $\mathbf{v}=(3,-4)$.
  \item For $f(x,y)=\ln(x^{2}+3y)$ at $P=(2,0)$, give the slope of the
        graph in the direction of steepest ascent and the tangent line to
        the graph in that direction.
  \item For $f(x,y)=(x+y)\,\mathrm{e}^{x-y}$ at $P=(1,1)$, give the slope in
        the direction from $P$ towards $Q=(3,2)$ and the tangent line in
        that direction, once with the unit vector and once with
        $\mathbf{h}=Q-P$. Show that both describe the same line.
  \item For $Q(K,L)=4K^{1/2}L^{1/4}$ at $(16,16)$, give the tangent line to
        the graph in the $L$-direction, read it as an output estimate for
        $L=17$, and compare with the exact value.
\end{problems}

In problems 15--17, for the point $P$: (a) find the level $c=f(P)$; (b)
solve the level curve $f=c$ explicitly; (c) compute $\grad f(P)$ and check
that it is perpendicular to the tangent of the level curve; (d) give the
tangent line to the level curve in normal and in parametric form; (e) give
$D_{\mathbf{u}}f(P)$ in the direction of the gradient.

\begin{problems}[resume]
  \item $f(x,y)=\sqrt{2x^{2}+3y+1}$, \;$P=(1,2)$.
  \item $f(x,y)=x^{1/3}y^{2/3}$ on $x,y>0$, \;$P=(1,8)$.
  \item $f(x,y)=x\,\mathrm{e}^{y}+y$, \;$P=(2,0)$; solve the level curve
        for $x$.
\end{problems}

Problems 18--20: differentiability of vector-valued maps.

\begin{problems}[resume]
  \item Explain why
        $\mathbf{F}(x,y)=\bigl(x^{2}y,\;x-3y^{2},\;xy\bigr)$ is
        differentiable at every point. Give $J_{\mathbf{F}}(1,1)$, and
        compute the remainder vector $\mathbf{r}(\mathbf{h})$ at $(1,1)$
        to show directly that $\|\mathbf{r}(\mathbf{h})\|/\|\mathbf{h}\|\to0$.
  \item For $\mathbf{F}(x,y)=\bigl(\ln(3x^{2}+y^{4}),\;xy^{2},\;|y-1|\bigr)$,
        give the domain, the set on which $\mathbf{F}$ is differentiable,
        and the Jacobian there.
  \item Answer in words: when is a map $\mathbf{f}\colon\R^{n}\to\R^{m}$
        differentiable at $\mathbf{x}_{0}$; what is its derivative called
        and what shape does it have; how is it related to the partial
        derivatives of the components; and what is the geometric meaning
        for $n=m=1$ and for $n=2$, $m=1$?
\end{problems}

Problems 21--26: Jacobians of all shapes.

\begin{problems}[resume]
  \item $\mathbf{r}(t)=\bigl(\ln(1+t^{2}),\;\mathrm{e}^{-2t},\;\cos3t,\;5\bigr)$:
        the Jacobian, its shape, and its value at $t=0$.
  \item $\mathbf{F}(x,y,z)=(xyz,\;x^{2}-z)$ at $(1,2,3)$, and
        $g(x_{1},x_{2},x_{3},x_{4})=x_{1}x_{2}+x_{3}\mathrm{e}^{x_{4}}$ at
        $(1,1,1,0)$: shapes and values of the Jacobians.
  \item $\mathbf{F}(s,t)=\bigl(s+t,\;st,\;s^{2},\;\mathrm{e}^{s-t}\bigr)$:
        the Jacobian, its shape, and its value at $(1,1)$.
  \item Jacobians in $n$ dimensions: (a) the identity map
        $\mathbf{F}(\mathbf{x})=\mathbf{x}$ on $\R^{n}$; (b) the map
        $\mathbf{F}\colon\R^{n}\to\R^{n}$ with components
        $F_{k}(\mathbf{x})=k\,\mathrm{e}^{x_{k}}$, that is
        $(\mathrm{e}^{x_{1}},2\mathrm{e}^{x_{2}},\ldots,n\,\mathrm{e}^{x_{n}})$;
        (c) the curve $\mathbf{r}\colon\R\to\R^{n}$ with components
        $r_{k}(t)=t^{k}/k$, that is $(t,t^{2}/2,\ldots,t^{n}/n)$; (d)
        $f(\mathbf{x})=\sum_{k=1}^{n}x_{k}^{k}/k$.
  \item Jacobians of linear and affine maps: (a) $\mathbf{F}(\mathbf{x})=
        A\mathbf{x}$ with $A=\begin{pmatrix}2&-1&0\\1&0&3\\0&4&-2
        \end{pmatrix}$; (b) $\mathbf{F}(\mathbf{x})=3\mathbf{x}-2\mathbf{a}$
        for a fixed $\mathbf{a}\in\R^{n}$; (c) $f(\mathbf{x})=
        \mathbf{a}^{T}\mathbf{x}$; (d) $\mathbf{G}(\mathbf{x})=B\mathbf{x}
        +\mathbf{c}$ with $B$ a $2\times3$ matrix.
  \item Let $\mathbf{F}\colon\R^{3}\to\R^{2}$, $\mathbf{G}\colon\R^{2}\to
        \R^{4}$, $H\colon\R^{4}\to\R$ and $\mathbf{K}\colon\R\to\R^{3}$.
        Which of $\mathbf{G}\circ\mathbf{F}$, $\mathbf{F}\circ\mathbf{G}$,
        $H\circ\mathbf{G}$, $\mathbf{G}\circ H$, $\mathbf{K}\circ H$,
        $H\circ\mathbf{K}$, $\mathbf{F}\circ\mathbf{K}$,
        $\mathbf{K}\circ\mathbf{F}$ and $H\circ\mathbf{G}\circ\mathbf{F}$
        are defined, and what is the shape of the Jacobian of each one that
        is?
\end{problems}

Problems 27--30: the matrix chain rule with rectangular Jacobians. Verify
each result by composing first.

\begin{problems}[resume]
  \item $\mathbf{f}(x,y,z)=(x+yz,\;xz)$ and
        $\mathbf{g}(u,v)=(u,\;v,\;uv,\;u-v)$: the Jacobian of
        $\mathbf{g}\circ\mathbf{f}$ at $(1,2,1)$.
  \item $g(x,y,z)=x+2y+z^{2}$ and $\mathbf{f}(s)=(s,\;s^{2},\;\mathrm{e}^{s})$:
        say why both $\mathbf{f}\circ g$ and $g\circ\mathbf{f}$ exist; give
        the Jacobian of $\mathbf{f}\circ g$ at $(1,0,0)$ and the derivative
        of $g\circ\mathbf{f}$.
  \item $\mathbf{f}(x,y)=(x,\;y,\;xy)$ and $\mathbf{g}(u,v,w)=(u+w,\;vw)$:
        the Jacobian of $\mathbf{g}\circ\mathbf{f}$ at $(1,2)$ and of
        $\mathbf{f}\circ\mathbf{g}$ at $(1,1,1)$.
  \item $\mathbf{r}(t)=(\cos t,\;\sin t,\;t)$ and $F(x,y,z)=x^{2}+y^{2}+z^{2}$:
        $\dd{}{t}F(\mathbf{r}(t))$ by the chain rule, simplified.
\end{problems}

In problems 31--33, compute $D_{\mathbf{u}}f(P)$ and $D^{2}_{\mathbf{u}}f(P)$
for the direction of $\mathbf{h}$, both with the gradient and Hessian and
with the profile $\varphi(t)=f(P+t\mathbf{h})$, and say whether $f$ is
convex or concave in that direction at $P$.

\begin{problems}[resume]
  \item $f(x,y,z)=x\ln y+yz^{2}$, \;$P=(2,1,1)$, \;$\mathbf{h}=(1,2,-2)$.
  \item $f(x,y,z)=x\,\mathrm{e}^{yz}$, \;$P=(1,0,2)$, \;$\mathbf{h}=(2,1,2)$.
  \item $f(x,y,z)=xyz+\ln(xz)$, \;$P=(1,1,1)$, \;$\mathbf{h}=(1,-1,1)$.
  \item Redo practice problem 15 ($f=x^{2}y$ at $(1,2)$) with the
        \emph{unnormalised} $\mathbf{h}=(3,4)$: compute
        $\grad f\cdot\mathbf{h}$, $\mathbf{h}^{T}H\mathbf{h}$, and the
        profile derivatives $\varphi'(0)$, $\varphi''(0)$. Explain the
        factors $5$ and $25$.
\end{problems}

In problems 35--37, give the directional verdicts asked for, then decide
whether $f$ is convex, concave or neither on its whole domain.

\begin{problems}[resume]
  \item $f(x,y)=x^{3}-3x+y^{2}$: at $Q=(-1,1)$ in the direction of
        $\mathbf{h}=(3,4)$, compute $D_{\mathbf{u}}f$ and
        $D^{2}_{\mathbf{u}}f$; then use the point $(1,0)$.
  \item $f(x,y)=x^{4}-4xy+y^{2}$: classify the critical points, then give
        $D^{2}_{\mathbf{u}}f(0,0)$ for $\mathbf{u}=(1,1)/\sqrt2$ and for
        $\mathbf{u}=(0,1)$.
  \item $f(x,y)=x\ln y$ on $y>0$, at $P=(1,1)$: $D_{\mathbf{u}}f$ and
        $D^{2}_{\mathbf{u}}f$ for $\mathbf{u}=(1,0)$, $(0,1)$ and
        $(1,1)/\sqrt2$.
\end{problems}

Problems 38--40: statements to check. Justify every answer.

\begin{problems}[resume]
  \item True or false? (a) Walking along the contour line through $P$,
        the rate of change of $f$ is $0$. (b) The direction of steepest
        descent is $-\grad f(P)$. (c) If $\grad f(P)=\mathbf{0}$, then $P$
        is a peak or a valley. (d) $|\grad f(P)|$ is the slope of the
        steepest path through $P$. (e) The gradient at $P$ is tangent to
        the level curve through $P$.
  \item A hiker stands on a hillside with neither map nor formula. How can
        she find the direction of the gradient, and how does the slope in
        other directions compare with the slope in that one?
  \item True or false? (a) For $f(x,y)$ of two variables: if
        $D_{\mathbf{u}}f(P)=0$ for two non-parallel unit vectors $\mathbf{u}$, then $\grad f(P)=\mathbf{0}$. (b) If
        $D_{\mathbf{u}}f(P)=3$, then $D_{-\mathbf{u}}f(P)=-3$. (c) A
        function can have $D_{\mathbf{u}}f(P)=5$ for some unit vector
        $\mathbf{u}$ while $|\grad f(P)|=4$.
\end{problems}

Problems 41--42: an input change vector of fixed length is to raise output
as fast as possible.

\begin{problems}[resume]
  \item $Y=10K^{0.4}L^{0.6}$ at $K=32768$, $L=7776$ (so that
        $K^{0.4}=64$, $L^{0.6}=216$). The change vector
        $(\mathrm{d}K,\mathrm{d}L)$ has length $100$. Find the best ratio
        $\mathrm{d}K:\mathrm{d}L$, the amounts, the estimated and the exact
        rise in output; then compare with spending a budget
        $\mathrm{d}K+\mathrm{d}L=100$.
  \item $Y=\sqrt{KL}$ at $K=400$, $L=900$, with a change vector of length
        $20$: the same questions.
\end{problems}

\section{Additional Solutions}

\begin{solutions}
  \item $\grad f=(2x+3y,\,-2y+3x)$; at $(1,2)$:
        $\grad f=(2+6,\,-4+3)=(8,-1)$. Normalise:
        $|\mathbf{v}|=\sqrt{16+9}=5$, $\mathbf{u}=(4/5,\,3/5)$.
        \[
          D_{\mathbf{u}}f=8\cdot\frac{4}{5}+(-1)\cdot\frac{3}{5}
            =\frac{32}{5}-\frac{3}{5}=\frac{29}{5} .
        \]

  \item $\grad f=(2x,\,4y,\,6z)$; at $(1,1,1)$: $\grad f=(2,4,6)$.
        Normalise: $|\mathbf{v}|=\sqrt{1+1+1}=\sqrt{3}$,
        $\mathbf{u}=(1,-1,1)/\sqrt{3}$.
        \[
          D_{\mathbf{u}}f=\frac{2-4+6}{\sqrt{3}}
            =\frac{4}{\sqrt{3}}=\frac{4\sqrt{3}}{3} .
        \]

  \item $\grad U=(1/x,\,2/y)$; at $(1,2)$: $\grad U=(1,1)$,
        $|\grad U|=\sqrt{2}$. Steepest ascent along
        \[
          \mathbf{u}=\frac{(1,1)}{\sqrt{2}},\qquad
          \text{maximum rate}=\sqrt{2} .
        \]

  \item $\grad f=(y,\,x)$; at $(2,1)$: $\grad f=(1,2)$. We need
        $\mathbf{u}\perp(1,2)$: $u_{1}+2u_{2}=0$, so
        $\mathbf{u}=\pm(2,-1)/\sqrt{5}$.
        \textbf{Check.} $(1,2)\cdot(2,-1)=2-2=0$. \checkmark

  \item The directional derivatives toward $(1,0)$ and $(0,1)$ are
        precisely $f_{x}(P)$ and $f_{y}(P)$, so $\grad f(P)=(3,-4)$.
        Maximum rate: $|\grad f(P)|=\sqrt{9+16}=5$, attained toward
        $(3,-4)/5$. Toward $\mathbf{u}=(3/5,\,4/5)$:
        \[
          D_{\mathbf{u}}f=3\cdot\frac{3}{5}+(-4)\cdot\frac{4}{5}
            =\frac{9}{5}-\frac{16}{5}=-\frac{7}{5} .
        \]

  \item $\grad f=\Bigl(\dfrac{2x}{x^{2}+y^{2}},\,
        \dfrac{2y}{x^{2}+y^{2}}\Bigr)$; at $(1,0)$: $\grad f=(2,0)$, so
        \[
          D_{\mathbf{u}}f=(2,0)\cdot(0,1)=0 :
        \]
        $\mathbf{u}$ is tangent to the level curve $x^{2}+y^{2}=1$
        through $(1,0)$. Second partials:
        \[
          f_{xx}=\frac{2y^{2}-2x^{2}}{(x^{2}+y^{2})^{2}},\qquad
          f_{yy}=\frac{2x^{2}-2y^{2}}{(x^{2}+y^{2})^{2}},\qquad
          f_{xy}=\frac{-4xy}{(x^{2}+y^{2})^{2}} ;
        \]
        at $(1,0)$: $f_{xx}=-2$, $f_{yy}=2$, $f_{xy}=0$. Hence with
        $\mathbf{u}=(0,1)$:
        \[
          D^{2}_{\mathbf{u}}f=f_{yy}=2>0 .
        \]
        Along $\mathbf{u}$ the function is flat to first order but
        curves upward: moving off $(1,0)$ vertically raises $f$ at
        second order. \textbf{Check.}
        $f(1,t)=\ln(1+t^{2})\approx t^{2}$ for small $t$: derivative
        $0$, second derivative $2$. \checkmark

  \item
        \[
          J_{\mathbf{f}}=\begin{pmatrix}2x&-2y\\3y&3x+2y\end{pmatrix},
          \qquad
          J_{\mathbf{f}}(1,1)=\begin{pmatrix}2&-2\\3&5\end{pmatrix} .
        \]

  \item $J_{f}=\bigl(y\;\;x\bigr)$, evaluated at
        $\mathbf{r}(t)=(t^{2},\,t^{3})$: $\bigl(t^{3}\;\;t^{2}\bigr)$;
        $J_{\mathbf{r}}(t)=\begin{pmatrix}2t\\3t^{2}\end{pmatrix}$.
        \[
          \dd{z}{t}=\bigl(t^{3}\;\;t^{2}\bigr)
            \begin{pmatrix}2t\\3t^{2}\end{pmatrix}
            =2t^{4}+3t^{4}=5t^{4} .
        \]
        \textbf{Check.} $z=t^{2}\cdot t^{3}=t^{5}$, so
        $\dd{z}{t}=5t^{4}$. \checkmark

  \item Inner and outer Jacobians:
        \[
          J_{\mathbf{g}}(s,t)=\begin{pmatrix}t&s\\1&1\end{pmatrix},\qquad
          J_{\mathbf{F}}(x,y)=\begin{pmatrix}1&1\\y&x\end{pmatrix} .
        \]
        Evaluate the outer at $(x,y)=\mathbf{g}(s,t)=(st,\,s+t)$ and
        multiply:
        \[
          J_{\mathbf{F}\circ\mathbf{g}}
            =\begin{pmatrix}1&1\\s+t&st\end{pmatrix}
             \begin{pmatrix}t&s\\1&1\end{pmatrix}
            =\begin{pmatrix}t+1&s+1\\
              (s+t)t+st&(s+t)s+st\end{pmatrix}
            =\begin{pmatrix}t+1&s+1\\
              t^{2}+2st&s^{2}+2st\end{pmatrix} .
        \]
        \textbf{Check.} Composing first:
        $(\mathbf{F}\circ\mathbf{g})(s,t)
        =\bigl(st+s+t,\;s^{2}t+st^{2}\bigr)$, with Jacobian
        \[
          \begin{pmatrix}t+1&s+1\\2st+t^{2}&s^{2}+2st\end{pmatrix} .
          \quad\checkmark
        \]

  \item Inner map $\mathbf{g}(K,L)=(2K+L,\;K+3L)$ with constant
        Jacobian
        $J_{\mathbf{g}}=\begin{pmatrix}2&1\\1&3\end{pmatrix}$; outer
        function $y=x_{1}x_{2}$ with row Jacobian
        $\grad y=(x_{2}\;\;x_{1})$. At $(K,L)=(1,1)$:
        $(x_{1},x_{2})=(3,4)$, so the row is $(4\;\;3)$ and
        \[
          \Bigl(\pd{y}{K},\,\pd{y}{L}\Bigr)
            =(4\;\;3)\begin{pmatrix}2&1\\1&3\end{pmatrix}
            =\bigl(4\cdot 2+3\cdot 1,\;4\cdot 1+3\cdot 3\bigr)
            =(11,\;13) .
        \]
        \textbf{Check.} Expanding,
        $y=(2K+L)(K+3L)=2K^{2}+7KL+3L^{2}$, so
        $\pd{y}{K}=4K+7L=11$ and $\pd{y}{L}=7K+6L=13$ at $(1,1)$.
        \checkmark

  \item $f(P)=2+8=10$ and $\grad f=(2xy,\;x^{2}+3y^{2})=(4,13)$ at $P$.
        \begin{itemize}[leftmargin=1.2em,itemsep=0.2ex]
          \item $x$-direction: slope $f_{x}(P)=4$, line
                $(1+t,\;2,\;10+4t)$.
          \item $y$-direction: slope $f_{y}(P)=13$, line
                $(1,\;2+t,\;10+13t)$.
          \item $|\mathbf{v}|=5$, $\mathbf{u}=(\tfrac35,-\tfrac45)$,
                $D_{\mathbf{u}}f=\tfrac{12}{5}-\tfrac{52}{5}=-8$: the
                surface falls with slope $8$ that way. Line
                $\bigl(1+\tfrac35t,\;2-\tfrac45t,\;10-8t\bigr)$.
        \end{itemize}
  \item $f(P)=\ln4$. With $w=x^{2}+3y$, $\grad f=\bigl(\tfrac{2x}{w},\,
        \tfrac{3}{w}\bigr)=(1,\tfrac34)$ at $P$, where $w=4$. Steepest
        ascent: $\mathbf{u}=\grad f/|\grad f|$ with $|\grad f|=\tfrac54$, so
        $\mathbf{u}=(\tfrac45,\tfrac35)$ and the slope is
        $D_{\mathbf{u}}f=|\grad f|=\tfrac54$. Tangent line:
        \[
          (x,y,z)=\bigl(2+\tfrac45t,\;\tfrac35t,\;\ln4+\tfrac54t\bigr) .
        \]
  \item $f(P)=2\mathrm{e}^{0}=2$. Product rule:
        $f_{x}=\mathrm{e}^{x-y}(x+y+1)=3$,
        $f_{y}=\mathrm{e}^{x-y}(1-x-y)=-1$ at $P$. $\mathbf{h}=Q-P=(2,1)$,
        $|\mathbf{h}|=\sqrt5$, $\mathbf{u}=(2,1)/\sqrt5$:
        \[
          D_{\mathbf{u}}f=\frac{3\cdot2-1\cdot1}{\sqrt5}=\sqrt5\approx2.236 .
        \]
        With the unit vector:
        $\bigl(1+\tfrac{2}{\sqrt5}t,\;1+\tfrac{1}{\sqrt5}t,\;2+\sqrt5\,t\bigr)$.
        With $\mathbf{h}$, $\grad f\cdot\mathbf{h}=5$:
        $(1+2s,\;1+s,\;2+5s)$. Putting $s=t/\sqrt5$ in the second gives the
        first, so it is the same line; only the parameter runs at a
        different speed.
  \item $Q(16,16)=4\cdot4\cdot2=32$,
        $Q_{L}=K^{1/2}L^{-3/4}=4\cdot\tfrac18=\tfrac12$. Tangent line in the
        $L$-direction: $(16,\;16+t,\;32+\tfrac12t)$. At $L=17$ ($t=1$) it
        gives $32.5$; exactly, $Q(16,17)=16\cdot17^{1/4}\approx32.4887$. The
        slope $\tfrac12$ is the marginal product of labour, and the line
        slightly overestimates because $Q$ is concave in $L$.
  \item (a) $c=\sqrt{2+6+1}=3$. (b) $2x^{2}+3y+1=9$, so
        $y=\dfrac{8-2x^{2}}{3}$, a downward parabola. Its slope at $x=1$ is
        $y'=-\tfrac{4x}{3}=-\tfrac43$, tangent direction $(3,-4)$.
        (c) $\grad f=\Bigl(\dfrac{2x}{f},\,\dfrac{3}{2f}\Bigr)=
        \bigl(\tfrac23,\tfrac12\bigr)$, and
        $\bigl(\tfrac23,\tfrac12\bigr)\cdot(3,-4)=2-2=0$. \checkmark\ (d)
        Normal form: $\tfrac23(x-1)+\tfrac12(y-2)=0$, i.e.\
        $4(x-1)+3(y-2)=0$, $y=2-\tfrac43(x-1)$. Parametric:
        $(x,y)=(1,2)+t(3,-4)$. (e) $D_{\mathbf{u}}f=|\grad f|=
        \sqrt{\tfrac49+\tfrac14}=\sqrt{\tfrac{25}{36}}=\tfrac56$.
  \item (a) $c=1\cdot8^{2/3}=4$. (b) $y^{2/3}=4x^{-1/3}$, so
        $y=8x^{-1/2}$; slope $y'=-4x^{-3/2}=-4$ at $x=1$, direction $(1,-4)$.
        (c) $f_{x}=\tfrac13\dfrac fx=\tfrac43$,
        $f_{y}=\tfrac23\dfrac fy=\tfrac23\cdot\tfrac48=\tfrac13$, and
        $\bigl(\tfrac43,\tfrac13\bigr)\cdot(1,-4)=0$. \checkmark\ (d)
        $\tfrac43(x-1)+\tfrac13(y-8)=0$, i.e.\ $y=12-4x$; parametric
        $(1,8)+t(1,-4)$. (e) $|\grad f|=\sqrt{\tfrac{16}{9}+\tfrac19}=
        \tfrac{\sqrt{17}}{3}\approx1.374$.
  \item (a) $c=2\cdot1+0=2$. (b) $x\mathrm{e}^{y}=2-y$, so
        $x=(2-y)\mathrm{e}^{-y}$. Then
        $\dd{x}{y}=-\mathrm{e}^{-y}-(2-y)\mathrm{e}^{-y}=(y-3)\mathrm{e}^{-y}
        =-3$ at $y=0$: tangent direction $(\mathrm{d}x,\mathrm{d}y)=(-3,1)$.
        (c) $\grad f=(\mathrm{e}^{y},\;x\mathrm{e}^{y}+1)=(1,3)$, and
        $(1,3)\cdot(-3,1)=0$. \checkmark\ (d) $(x-2)+3y=0$, i.e.\
        $x=2-3y$; parametric $(2,0)+t(-3,1)$. (e) $|\grad f|=\sqrt{10}
        \approx3.162$.
  \item Each component is a polynomial, so its partials exist and are
        continuous everywhere; a $C^{1}$ map is differentiable, and a map is
        differentiable when each component is.
        $J_{\mathbf{F}}=\begin{pmatrix}2xy&x^{2}\\1&-6y\\y&x\end{pmatrix}$,
        and $J_{\mathbf{F}}(1,1)=\begin{pmatrix}2&1\\1&-6\\1&1\end{pmatrix}$.
        Remainder at $(1,1)$, component by component, with
        $\mathbf{F}(1,1)=(1,-2,1)$:
        \[\begin{aligned}
          r_{1}&=(1+h_{1})^{2}(1+h_{2})-1-2h_{1}-h_{2}
               =2h_{1}h_{2}+h_{1}^{2}+h_{1}^{2}h_{2},\\
          r_{2}&=(1+h_{1})-3(1+h_{2})^{2}+2-h_{1}+6h_{2}=-3h_{2}^{2},\\
          r_{3}&=(1+h_{1})(1+h_{2})-1-h_{1}-h_{2}=h_{1}h_{2} .
        \end{aligned}\]
        Each is bounded by a constant times $\|\mathbf{h}\|^{2}$ (for
        $\|\mathbf{h}\|\le1$), so
        $\|\mathbf{r}\|\le C\|\mathbf{h}\|^{2}$ and
        $\|\mathbf{r}\|/\|\mathbf{h}\|\to0$.
  \item \emph{Domain.} $3x^{2}+y^{4}>0$ fails only at the origin: $D=\R^{2}
        \setminus\{(0,0)\}$. \emph{Differentiability.} The first two
        components are $C^{1}$ on $D$. The third, $|y-1|$, has
        $\partial/\partial y$ equal to $+1$ above and $-1$ below the line
        $y=1$, and on the line the quotient $|h|/h$ has no limit. So
        $\mathbf{F}$ is differentiable exactly on
        $\{(x,y)\neq(0,0):y\neq1\}$, with
        \[
          J_{\mathbf{F}}=\begin{pmatrix}
            \dfrac{6x}{3x^{2}+y^{4}}&\dfrac{4y^{3}}{3x^{2}+y^{4}}\\[1.2ex]
            y^{2}&2xy\\ 0&\operatorname{sign}(y-1)\end{pmatrix} .
        \]
  \item $\mathbf{f}$ is differentiable at $\mathbf{x}_{0}$ if
        $\mathbf{f}(\mathbf{x}_{0}+\mathbf{h})=\mathbf{f}(\mathbf{x}_{0})
        +J\mathbf{h}+\mathbf{r}(\mathbf{h})$ for a matrix $J$, with
        $\|\mathbf{r}(\mathbf{h})\|/\|\mathbf{h}\|\to0$: near
        $\mathbf{x}_{0}$ the map is a linear map plus an error that vanishes
        faster than the step. The matrix is the derivative (total
        derivative), the Jacobian matrix, of shape $m\times n$. Its entries
        are the partials $\partial f_{i}/\partial x_{j}$; their existence is
        necessary but not sufficient, while continuity of all of them is
        sufficient, and $\mathbf{f}$ is differentiable exactly when each
        component is. For $n=m=1$ the linear part is the tangent line of
        the graph, for $n=2$, $m=1$ the tangent plane.
  \item $J_{\mathbf{r}}(t)=\Bigl(\dfrac{2t}{1+t^{2}},\;-2\mathrm{e}^{-2t},
        \;-3\sin3t,\;0\Bigr)^{T}$, a $4\times1$ column (four outputs, one
        input); the constant component gives $0$. At $t=0$:
        $(0,-2,0,0)^{T}$.
  \item $J_{\mathbf{F}}=\begin{pmatrix}yz&xz&xy\\2x&0&-1\end{pmatrix}$,
        $2\times3$; at $(1,2,3)$: $\begin{pmatrix}6&3&2\\2&0&-1\end{pmatrix}$.
        $J_{g}=\bigl(x_{2}\ \ x_{1}\ \ \mathrm{e}^{x_{4}}\ \
        x_{3}\mathrm{e}^{x_{4}}\bigr)$, a $1\times4$ row; at $(1,1,1,0)$:
        $(1\ \ 1\ \ 1\ \ 1)$.
  \item $J_{\mathbf{F}}=\begin{pmatrix}1&1\\t&s\\2s&0\\
        \mathrm{e}^{s-t}&-\mathrm{e}^{s-t}\end{pmatrix}$, $4\times2$ (four
        outputs, two inputs); at $(1,1)$:
        $\begin{pmatrix}1&1\\1&1\\2&0\\1&-1\end{pmatrix}$.
  \item (a) $F_{k}=x_{k}$, so $\partial F_{k}/\partial x_{j}$ is $1$ for
        $j=k$ and $0$ otherwise: $J=I_{n}$. (b) $F_{k}=k\mathrm{e}^{x_{k}}$
        depends on $x_{k}$ alone: $J=\operatorname{diag}\bigl(\mathrm{e}^{x_{1}},
        2\mathrm{e}^{x_{2}},\ldots,n\,\mathrm{e}^{x_{n}}\bigr)$, $n\times n$.
        (c) $r_{k}=t^{k}/k$ has $r_{k}'=t^{k-1}$: the $n\times1$ column
        $(1,t,t^{2},\ldots,t^{n-1})^{T}$. (d) The term with $x_{k}$ is
        $x_{k}^{k}/k$, with derivative $x_{k}^{k-1}$: the $1\times n$ row
        $(1\ \ x_{2}\ \ x_{3}^{2}\ \ \cdots\ \ x_{n}^{n-1})$.
  \item (a) Component $i$ is $\sum_{j}a_{ij}x_{j}$, with partials $a_{ij}$:
        $J=A$ at every point. (b) Component $k$ is $3x_{k}-2a_{k}$:
        $J=3I_{n}$, and the constant shift drops out. (c) $\partial f/
        \partial x_{k}=a_{k}$: $J=\mathbf{a}^{T}$, a $1\times n$ row. (d)
        $J=B$, $2\times3$, whatever $\mathbf{c}$ is.
  \item Match outputs of the inner map with inputs of the outer map.
        \begin{itemize}[leftmargin=1.2em,itemsep=0.2ex]
          \item $\mathbf{G}\circ\mathbf{F}$: $2=2$, defined,
                $\R^{3}\to\R^{4}$, $4\times3$. $\mathbf{F}\circ\mathbf{G}$:
                $4\neq3$, undefined.
          \item $H\circ\mathbf{G}$: $4=4$, defined, $\R^{2}\to\R$,
                $1\times2$. $\mathbf{G}\circ H$: $1\neq2$, undefined.
          \item $\mathbf{K}\circ H$: $1=1$, defined, $\R^{4}\to\R^{3}$,
                $3\times4$. $H\circ\mathbf{K}$: $3\neq4$, undefined.
          \item $\mathbf{F}\circ\mathbf{K}$: $3=3$, defined, $\R\to\R^{2}$,
                $2\times1$. $\mathbf{K}\circ\mathbf{F}$: $2\neq1$,
                undefined.
          \item $H\circ\mathbf{G}\circ\mathbf{F}$: $\R^{3}\to\R^{2}\to\R^{4}
                \to\R$, defined, $1\times3=(1\times4)(4\times2)(2\times3)$.
        \end{itemize}
  \item $\mathbf{f}(1,2,1)=(1+2,\;1)=(3,1)$.
        $J_{\mathbf{f}}=\begin{pmatrix}1&z&y\\z&0&x\end{pmatrix}
        =\begin{pmatrix}1&1&2\\1&0&1\end{pmatrix}$ and
        $J_{\mathbf{g}}=\begin{pmatrix}1&0\\0&1\\v&u\\1&-1\end{pmatrix}
        =\begin{pmatrix}1&0\\0&1\\1&3\\1&-1\end{pmatrix}$ at $(u,v)=(3,1)$.
        \[
          J_{\mathbf{g}\circ\mathbf{f}}=
          \begin{pmatrix}1&0\\0&1\\1&3\\1&-1\end{pmatrix}
          \begin{pmatrix}1&1&2\\1&0&1\end{pmatrix}
          =\begin{pmatrix}1&1&2\\1&0&1\\4&1&5\\0&1&1\end{pmatrix},
        \]
        a $(4\times2)(2\times3)=4\times3$ product. \textbf{Check.}
        $\mathbf{g}\circ\mathbf{f}=\bigl(x+yz,\;xz,\;x^{2}z+xyz^{2},\;
        x+yz-xz\bigr)$; the third row is
        $(2xz+yz^{2},\;xz^{2},\;x^{2}+2xyz)=(4,1,5)$ and the fourth
        $(1-z,\;z,\;y-x)=(0,1,1)$ at $(1,2,1)$. \checkmark
  \item $g\colon\R^{3}\to\R$ and $\mathbf{f}\colon\R\to\R^{3}$: in
        $\mathbf{f}\circ g$ the one output of $g$ feeds the one input of
        $\mathbf{f}$, and in $g\circ\mathbf{f}$ the three outputs of
        $\mathbf{f}$ feed the three inputs of $g$. Both exist.
        \emph{$\mathbf{f}\circ g\colon\R^{3}\to\R^{3}$.} The chain rule
        gives a column times a row,
        $J=\mathbf{f}'(g(\mathbf{x}))\,\grad g(\mathbf{x})^{T}$, a
        rank-one $3\times3$ matrix. At $(1,0,0)$: $g=1$,
        $\mathbf{f}'(1)=(1,\;2,\;\mathrm{e})^{T}$, $\grad g=(1,2,2z)=(1,2,0)$:
        \[
          J=\begin{pmatrix}1\\2\\\mathrm{e}\end{pmatrix}(1\ \ 2\ \ 0)
           =\begin{pmatrix}1&2&0\\2&4&0\\\mathrm{e}&2\mathrm{e}&0\end{pmatrix} .
        \]
        Check: $\mathbf{f}(g)=(g,\;g^{2},\;\mathrm{e}^{g})$ has rows
        $\grad g$, $2g\grad g$, $\mathrm{e}^{g}\grad g$. \checkmark\
        \emph{$g\circ\mathbf{f}\colon\R\to\R$.}
        $\grad g(\mathbf{f}(s))\cdot\mathbf{f}'(s)=(1,\;2,\;2\mathrm{e}^{s})
        \cdot(1,\;2s,\;\mathrm{e}^{s})=1+4s+2\mathrm{e}^{2s}$. Check:
        $g(\mathbf{f}(s))=s+2s^{2}+\mathrm{e}^{2s}$. \checkmark
  \item \emph{$\mathbf{g}\circ\mathbf{f}\colon\R^{2}\to\R^{2}$} at $(1,2)$:
        $\mathbf{f}(1,2)=(1,2,2)$,
        $J_{\mathbf{f}}=\begin{pmatrix}1&0\\0&1\\y&x\end{pmatrix}
        =\begin{pmatrix}1&0\\0&1\\2&1\end{pmatrix}$,
        $J_{\mathbf{g}}=\begin{pmatrix}1&0&1\\0&w&v\end{pmatrix}
        =\begin{pmatrix}1&0&1\\0&2&2\end{pmatrix}$, so
        \[
          J_{\mathbf{g}\circ\mathbf{f}}=\begin{pmatrix}1&0&1\\0&2&2\end{pmatrix}
          \begin{pmatrix}1&0\\0&1\\2&1\end{pmatrix}
          =\begin{pmatrix}3&1\\4&4\end{pmatrix} .
        \]
        Check: $\mathbf{g}(\mathbf{f}(x,y))=(x+xy,\;xy^{2})$, Jacobian
        $\begin{pmatrix}1+y&x\\y^{2}&2xy\end{pmatrix}$, which is
        $\begin{pmatrix}3&1\\4&4\end{pmatrix}$ at $(1,2)$. \checkmark\
        \emph{$\mathbf{f}\circ\mathbf{g}\colon\R^{3}\to\R^{3}$} at
        $(1,1,1)$: $\mathbf{g}(1,1,1)=(2,1)$, so
        $J_{\mathbf{f}}(2,1)=\begin{pmatrix}1&0\\0&1\\1&2\end{pmatrix}$ and
        $J_{\mathbf{g}}(1,1,1)=\begin{pmatrix}1&0&1\\0&1&1\end{pmatrix}$:
        \[
          J_{\mathbf{f}\circ\mathbf{g}}=\begin{pmatrix}1&0\\0&1\\1&2\end{pmatrix}
          \begin{pmatrix}1&0&1\\0&1&1\end{pmatrix}
          =\begin{pmatrix}1&0&1\\0&1&1\\1&2&3\end{pmatrix} .
        \]
        Check: $\mathbf{f}(\mathbf{g})=\bigl(u+w,\;vw,\;(u+w)vw\bigr)$; the
        last row is $\bigl(vw,\;(u+w)w,\;vw+(u+w)v\bigr)=(1,2,3)$.
        \checkmark
  \item $\grad F=(2x,2y,2z)$ and $\mathbf{r}'(t)=(-\sin t,\cos t,1)$:
        \[
          \dd{}{t}F(\mathbf{r}(t))=2\cos t(-\sin t)+2\sin t\cos t+2t\cdot1
          =2t .
        \]
        Check: $F(\mathbf{r}(t))=\cos^{2}t+\sin^{2}t+t^{2}=1+t^{2}$.
        \checkmark
  \item $|\mathbf{h}|=3$. $\grad f=(\ln y,\;x/y+z^{2},\;2yz)=(0,3,2)$, so
        $\grad f\cdot\mathbf{h}=0+6-4=2$ and $D_{\mathbf{u}}f=\tfrac23$.
        $H=\begin{pmatrix}0&1/y&0\\1/y&-x/y^{2}&2z\\0&2z&2y\end{pmatrix}
        =\begin{pmatrix}0&1&0\\1&-2&2\\0&2&2\end{pmatrix}$ and
        \[
          \mathbf{h}^{T}H\mathbf{h}=2\cdot1\cdot(1)(2)-2\cdot4
          +2\cdot2\cdot(2)(-2)+2\cdot4=4-8-16+8=-12,
        \]
        so $D^{2}_{\mathbf{u}}f=-\tfrac{12}{9}=-\tfrac43$. \emph{Profile.}
        $\varphi(t)=(2+t)\ln(1+2t)+(1+2t)(1-2t)^{2}$, and
        $(1+2t)(1-2t)^{2}=1-2t-4t^{2}+8t^{3}$. The log part has derivative
        $\ln(1+2t)+\dfrac{2(2+t)}{1+2t}$, equal to $4$ at $0$, and second
        derivative $\dfrac{2}{1+2t}+\dfrac{2(1+2t)-4(2+t)}{(1+2t)^{2}}$,
        equal to $2-6=-4$. So $\varphi'(0)=4-2=2$ and
        $\varphi''(0)=-4-8=-12$. \checkmark\ Concave in this direction.
  \item $|\mathbf{h}|=3$. $\grad f=\mathrm{e}^{yz}(1,\;xz,\;xy)=(1,2,0)$;
        $D_{\mathbf{u}}f=\tfrac{2+2+0}{3}=\tfrac43$. Second partials at $P$
        ($\mathrm{e}^{yz}=1$): $f_{xx}=0$, $f_{xy}=z\mathrm{e}^{yz}=2$,
        $f_{xz}=y\mathrm{e}^{yz}=0$, $f_{yy}=xz^{2}\mathrm{e}^{yz}=4$,
        $f_{yz}=x(1+yz)\mathrm{e}^{yz}=1$, $f_{zz}=xy^{2}\mathrm{e}^{yz}=0$.
        \[
          \mathbf{h}^{T}H\mathbf{h}=2\cdot2\cdot(2)(1)+4\cdot1+2\cdot1\cdot(1)(2)
          =8+4+4=16,\qquad D^{2}_{\mathbf{u}}f=\tfrac{16}{9} .
        \]
        \emph{Profile.} $\varphi(t)=(1+2t)\,\mathrm{e}^{t(2+2t)}
        =(1+2t)\,\mathrm{e}^{2t+2t^{2}}$. With $E=\mathrm{e}^{2t+2t^{2}}$,
        $E'=(2+4t)E$:
        $\varphi'=\bigl(2+(1+2t)(2+4t)\bigr)E$, so $\varphi'(0)=4$, and
        $\varphi''=\bigl(4+8t+4+8t\bigr)E+\bigl(2+(1+2t)(2+4t)\bigr)(2+4t)E$,
        so $\varphi''(0)=8+4\cdot2=16$. \checkmark\ Convex in this direction.
  \item $|\mathbf{h}|=\sqrt3$. $\grad f=(yz+\tfrac1x,\;xz,\;xy+\tfrac1z)
        =(2,1,2)$, $\grad f\cdot\mathbf{h}=2-1+2=3$,
        $D_{\mathbf{u}}f=3/\sqrt3=\sqrt3$.
        $H=\begin{pmatrix}-1/x^{2}&z&y\\z&0&x\\y&x&-1/z^{2}\end{pmatrix}
        =\begin{pmatrix}-1&1&1\\1&0&1\\1&1&-1\end{pmatrix}$, and
        \[
          \mathbf{h}^{T}H\mathbf{h}=-1+0-1+2\bigl[(1)(-1)+(1)(1)+(-1)(1)\bigr]
          =-2-2=-4,\qquad D^{2}_{\mathbf{u}}f=-\tfrac43 .
        \]
        \emph{Profile.} $\varphi(t)=(1+t)(1-t)(1+t)+\ln\bigl((1+t)^{2}\bigr)
        =1+t-t^{2}-t^{3}+2\ln(1+t)$. $\varphi'(t)=1-2t-3t^{2}+\tfrac{2}{1+t}$,
        $\varphi'(0)=3$; $\varphi''(t)=-2-6t-\tfrac{2}{(1+t)^{2}}$,
        $\varphi''(0)=-4$. \checkmark\ Concave in this direction.
  \item At $(1,2)$: $\grad f=(2xy,x^{2})=(4,1)$ and
        $H=\begin{pmatrix}2y&2x\\2x&0\end{pmatrix}=\begin{pmatrix}4&2\\2&0
        \end{pmatrix}$. With $\mathbf{h}=(3,4)$:
        $\grad f\cdot\mathbf{h}=12+4=16$ and
        $\mathbf{h}^{T}H\mathbf{h}=4\cdot9+2\cdot2\cdot12+0=84$. Profile:
        $\varphi(t)=(1+3t)^{2}(2+4t)=2+16t+42t^{2}+36t^{3}$, so
        $\varphi'(0)=16$ and $\varphi''(0)=84$. \checkmark\ Practice
        problem 15 found $D_{\mathbf{u}}f=\tfrac{16}{5}$ and
        $D^{2}_{\mathbf{u}}f=\tfrac{84}{25}$ with $\mathbf{u}=\mathbf{h}/5$:
        moving along $\mathbf{h}$ covers $5$ units of distance per unit of
        $t$, so first derivatives scale by $|\mathbf{h}|=5$ and second
        derivatives by $|\mathbf{h}|^{2}=25$. The sign, and so the
        convexity verdict, is the same either way.
  \item $\grad f=(3x^{2}-3,\;2y)$, $H=\begin{pmatrix}6x&0\\0&2\end{pmatrix}$.
        At $Q=(-1,1)$: $\grad f=(0,2)$, $|\mathbf{h}|=5$,
        $\mathbf{u}=(\tfrac35,\tfrac45)$, so $D_{\mathbf{u}}f=\tfrac85$;
        $H=\begin{pmatrix}-6&0\\0&2\end{pmatrix}$ and
        $D^{2}_{\mathbf{u}}f=\dfrac{-6\cdot9+2\cdot16}{25}=-\dfrac{22}{25}<0$:
        rising, but concave in this direction at $Q$. At $(1,0)$ the
        gradient is $\mathbf{0}$ and $H=\begin{pmatrix}6&0\\0&2\end{pmatrix}$
        is PD: a strict local minimum, with positive curvature in every
        direction. Negative curvature somewhere rules out convexity,
        positive curvature somewhere rules out concavity: $f$ is
        \textbf{neither}.
  \item $f_{x}=4x^{3}-4y$, $f_{y}=-4x+2y$. From $f_{y}=0$: $y=2x$; then
        $4x^{3}=8x$, so $x=0$ or $x^{2}=2$. Critical points $(0,0)$ and
        $\pm(\sqrt2,2\sqrt2)$. $H=\begin{pmatrix}12x^{2}&-4\\-4&2\end{pmatrix}$.
        At $\pm(\sqrt2,2\sqrt2)$: $\begin{pmatrix}24&-4\\-4&2\end{pmatrix}$,
        $D_{1}=24$, $D_{2}=48-16=32$: PD, strict local minima, with value
        $4-16+8=-4$. At $(0,0)$: $\begin{pmatrix}0&-4\\-4&2\end{pmatrix}$,
        $D_{2}=-16$: a saddle. There, for $\mathbf{u}=(1,1)/\sqrt2$,
        $D^{2}_{\mathbf{u}}f=\tfrac{0-8+2}{2}=-3<0$ (concave that way), and
        for $\mathbf{u}=(0,1)$, $D^{2}_{\mathbf{u}}f=2>0$ (convex that way).
        Globally: negative curvature at the origin and positive definite
        Hessians at the minima, so $f$ is neither convex nor concave.
  \item $\grad f=(\ln y,\;x/y)=(0,1)$ and
        $H=\begin{pmatrix}0&1/y\\1/y&-x/y^{2}\end{pmatrix}=
        \begin{pmatrix}0&1\\1&-1\end{pmatrix}$ at $P$.
        \begin{itemize}[leftmargin=1.2em,itemsep=0.2ex]
          \item $\mathbf{u}=(1,0)$: $D_{\mathbf{u}}f=0$,
                $D^{2}_{\mathbf{u}}f=0$. In fact $f(x,1)=0$: flat, no
                verdict beyond ``linear''.
          \item $\mathbf{u}=(0,1)$: $D_{\mathbf{u}}f=1$,
                $D^{2}_{\mathbf{u}}f=-1$: concave in this direction.
          \item $\mathbf{u}=(1,1)/\sqrt2$: $D_{\mathbf{u}}f=\tfrac{1}{\sqrt2}$,
                $D^{2}_{\mathbf{u}}f=\tfrac12(0+2-1)=\tfrac12$: convex in
                this direction.
        \end{itemize}
        Both signs at one point already decide it: $f$ is neither convex
        nor concave. Indeed $\det H=-1/y^{2}<0$ everywhere.
  \item (a) True: along the contour $f$ is constant; equivalently the
        tangent $\mathbf{u}$ is perpendicular to $\grad f$, so
        $D_{\mathbf{u}}f=0$. (b) True when $\grad f(P)\neq\mathbf{0}$: the
        unit vector $-\grad f/|\grad f|$ gives the most negative rate,
        $-|\grad f|$. (c) False: $f=xy$ has $\grad f(0,0)=\mathbf{0}$ at a
        saddle. (d) True: the largest directional derivative is
        $|\grad f(P)|$. (e) False: it is perpendicular (normal) to the level
        curve.
  \item She walks a short way in the direction in which she neither climbs
        nor descends: she is following the contour line. The gradient
        points at right angles to that path, on the uphill side. Its
        direction is the steepest one; in a direction at angle $\theta$ to
        it the slope is $|\grad f|\cos\theta$: smaller, zero along the
        contour, and negative once she turns downhill.
  \item (a) True: $\grad f\cdot\mathbf{u}_{1}=\grad f\cdot\mathbf{u}_{2}=0$
        with $\mathbf{u}_{1},\mathbf{u}_{2}$ independent forces $\grad f$ to
        be perpendicular to the whole plane, hence $\mathbf{0}$. With three
        variables two directions are not enough: $f(x,y,z)=z$ has
        $D_{\mathbf{e}_{1}}f=D_{\mathbf{e}_{2}}f=0$ but $\grad f=(0,0,1)$.
        (b) True:
        $D_{-\mathbf{u}}f=\grad f\cdot(-\mathbf{u})=-D_{\mathbf{u}}f$.
        (c) False: by Cauchy--Schwarz,
        $D_{\mathbf{u}}f=\grad f\cdot\mathbf{u}\le|\grad f|\,|\mathbf{u}|=4$.
  \item $Y=10\cdot64\cdot216=138240$. In output-ratio form
        $Y_{K}=0.4\,Y/K=\tfrac{27}{16}=1.6875$ and
        $Y_{L}=0.6\,Y/L=\tfrac{32}{3}\approx10.667$. Best ratio
        $\mathrm{d}K:\mathrm{d}L=Y_{K}:Y_{L}=81:512$.
        $|\grad Y|=\sqrt{1.6875^{2}+10.667^{2}}\approx10.7993$, so
        \[
          (\mathrm{d}K,\mathrm{d}L)=\frac{100}{10.7993}(1.6875,\;10.667)
          \approx(15.63,\;98.77),
          \qquad \mathrm{d}Y\approx100\cdot10.7993\approx1079.9 .
        \]
        Exactly, $Y(32783.63,\;7874.77)-Y\approx1077.5$. A budget
        $\mathrm{d}K+\mathrm{d}L=100$ is maximised in the linear estimate by
        spending everything on labour, the input with the larger marginal
        product: $\mathrm{d}Y\approx1066.7$ (exactly $1063.9$). That is
        less, but it also uses less: the steepest-ascent vector has
        $\mathrm{d}K+\mathrm{d}L\approx114.4$.
  \item $Y=\sqrt{400\cdot900}=600$, $Y_{K}=\dfrac{Y}{2K}=0.75$,
        $Y_{L}=\dfrac{Y}{2L}=\tfrac13$. Ratio $\mathrm{d}K:\mathrm{d}L=
        0.75:\tfrac13=9:4$. $|\grad Y|=\sqrt{0.5625+0.1111}\approx0.82074$,
        so $(\mathrm{d}K,\mathrm{d}L)\approx\tfrac{20}{0.82074}(0.75,\;0.3333)
        \approx(18.28,\;8.12)$ and $\mathrm{d}Y\approx20\cdot0.82074\approx
        16.41$; exactly $16.32$. A budget $\mathrm{d}K+\mathrm{d}L=20$ would
        go entirely into capital: $\mathrm{d}Y\approx15$ (exactly $14.82$).
\end{solutions}
